% ─────────────────────────────────────────────────────────────────
% Preambuła wspólna dla dokumentów projektu Feedback
% ─────────────────────────────────────────────────────────────────

% Kodowanie i język
\usepackage[utf8]{inputenc}
\usepackage[T1]{fontenc}
\usepackage[main=polish,provide=*]{babel}

% Układ strony
\usepackage{geometry}
\geometry{margin=2.5cm}

% Tabele
\usepackage{booktabs}
\usepackage{tabularray}
\usepackage{longtable}
\usepackage{array}
\usepackage{colortbl}  % dla \cellcolor w komórkach

% Grafika i strony specjalne
\usepackage{afterpage}
\usepackage{rotating}
\usepackage{pdflscape}

% Kolory
\usepackage{xcolor}
\definecolor{accent}{RGB}{30, 60, 100}
\definecolor{lowrisk}{RGB}{200, 230, 200}
\definecolor{medrisk}{RGB}{255, 235, 170}
\definecolor{highrisk}{RGB}{245, 195, 195}

% Tytuły, listy, akapity
\usepackage{titlesec}
\usepackage{parskip}
\usepackage{enumitem}
\usepackage{fancyhdr}

% Hyperref i cleveref – kolejność ma znaczenie:
% hyperref musi być załadowany przed cleveref
\usepackage{hyperref}
\hypersetup{plainpages=false, hypertexnames=false, colorlinks=true, linkcolor=accent, urlcolor=accent, citecolor=accent}
\usepackage[nameinlink]{cleveref}
% Polskie nazwy dla cleveref – formy skrócone unikają problemu z odmianą
\crefname{section}{sek.}{sek.}
\Crefname{section}{Sek.}{Sek.}
\crefname{subsection}{sek.}{sek.}
\Crefname{subsection}{Sek.}{Sek.}
\crefname{table}{tab.}{tab.}
\Crefname{table}{Tab.}{Tab.}
\crefname{figure}{rys.}{rys.}
\Crefname{figure}{Rys.}{Rys.}
\crefname{equation}{rów.}{rów.}
\Crefname{equation}{Rów.}{Rów.}
\crefname{enumi}{pkt}{pkt}
\Crefname{enumi}{Pkt}{Pkt}


% Numerowane sekcje z kolorem akcentu
\titleformat{\section}{\large\bfseries\color{accent}}{\thesection.}{0.6em}{}[\titlerule]
\titleformat{\subsection}{\normalsize\bfseries}{\thesubsection}{0.6em}{}
\titleformat{\subsubsection}{\normalsize\itshape}{\thesubsubsection}{0.6em}{}

% Łamanie wierszy – tolerancja dla polskich tekstów
\setlength{\emergencystretch}{3em}
\hyphenpenalty=1000
\tolerance=2000
